% !PRINTED 637
\weberpage{637}
%<<CONT
einer \gesperrt{Gauss}'schen Summe zu den Kreistheilungszahlen gehört,
und es ist
\weq{24}{\varrho^m = \Pi^{t}\, \mathfrak{p}^t_{t'}.}

Nach (9) und (11) ist aber, wenn darin $m$ durch $m_1$ ersetzt wird,
\weq{25}{\varrho_n^{m_1} = \pm\, p^{\frac{1}{2}\,m_1(m_2-1)+1}\, \psi_1(r_1^n)\,\dots\,\psi_{m_1-2}(r_1^n),}
\weq{26}{\varrho^n \varrho_n^{-1} = p^{\frac{1}{2}(n-1)(m_2-1)}\, \psi_1(r_1)\,\psi_2(r_1)\,\dots\,\psi_{n-1}(r_1)}
und es sind also $\varrho^n \varrho_n^{-1}$, $\varrho_n^{m_1}$ und folglich auch $\varrho_n^m$ Zahlen des
Körpers $\Omega_m$, und nach (25) bleibt die in (26) vorkommende
Verbindung $\varrho^n \varrho_n^{-1}$ in $\Omega_m$ enthalten, wenn $n$ um ein Vielfaches
von $m_1$ vermehrt wird. Es ist demnach $\varrho^n \varrho_n^{-1}$ auch dann noch
eine Zahl in $\Omega_m$, wenn $n$ nicht gerade die beschränkte Bedeutung
hat wie bisher, sondern eine beliebige durch $q$ nicht theilbare
Zahl bedeutet.

Diese Betrachtung hat den Zweck, das \gesperrt{Kummer}'sche
Theorem (16) gleichzeitig auf mehrere verschiedene Primzahlen
anwendbar zu machen, die zu verschiedenen Werthen von $m_1$
gehören, und das Resultat spricht sich dann in folgendem
Theorem aus:

3. \gesperrt{Es sei $\varphi_n$ ein System conjugirter Functionale des
Körpers $\Omega_m$, das mit keinen anderen natürlichen
Primzahlen verwandt ist, als solchen, die nach
dem Modul $q$ (oder bei $q = 2$ nach dem Modul 4)
mit 1 congruent sind; es durchlaufe ferner $t$
die durch $q$ nicht theilbaren Zahlen der Reihe
$1, 2, \dots, m - 1$, und $t'$ sei aus der Congruenz
$t\,t' \equiv 1 \pmod{m}$ bestimmt.}

\gesperrt{Dann ist das Functional des Körpers $\Omega_m$:}
\[ \text{a)}\qquad \Phi_n = \Pi^{t}\, \varphi^t_{nt'} = \varepsilon\,\vartheta_n^m \]
\gesperrt{mit der $m^{\text{ten}}$ Potenz einer Kreistheilungszahl $\vartheta_n$
associirt, die zwar selbst in einem höheren
Körper enthalten ist, deren $m^{\text{te}}$ Potenz aber dem
Körper $\Omega_m$ angehört, und es ist}
\[ \text{b)}\qquad \vartheta_n^{-1}\vartheta_1^n = \alpha \]
\gesperrt{eine Zahl des Körpers $\Omega_m$.}

Der Beweis ist in den vorangegangenen Ausführungen enthalten
und ergiebt sich aus den Formeln (25), (26), wenn man
%CONT>>
